\section{证明完成：费马大定理的推导}
\label{sec:ch10-proof}
% ============================================================================

\providecommand{\rhobar}{\bar{\rho}}

现在到了全书最庄重的一节。
前面各章分别证明了 Hecke 理论、模曲线几何、Eichler--Shimura 关系、Galois 表示与 Ribet 定理；
本章前半又补上了 Frey 曲线与 Wiles 的半稳定模性定理。
这些结果单独看都很美，
但它们真正令人屏息的一刻，
是在这里被无缝拼接成一条从假设到矛盾的闭链。

\textbf{我们遇到了什么问题？}
要证明费马大定理，
我们必须处理所有指数 $n\ge 3$。
这听起来像无穷多个方程，
但首先可以把它压缩成两个基本情形：
$n=4$ 与 $n=p\ge 5$ 为奇素数。
前者是 Fermat 本人用无穷递降法解决的；
后者则是现代模性机器真正介入的地方。
所以证明的最后结构其实非常清楚：
先做一次指数归约，
再把奇素数情形送入 Frey--Wiles--Ribet 链条。

\textbf{核心思想是什么？}
假设存在费马方程的非平凡解。
若指数是奇素数 $p$，
则这组解制造出一条 Frey 曲线 $E_{a,b,c}$。
由 \cref{thm:wiles}，它必须模；
由 \cref{thm:ribet}，它的 mod $p$ 表示可以把每个 $q\mid abc$ 从级别中删去；
最后只剩级别 $2$。
但第5章的维数公式告诉我们
\[
  \dim \Sk_2(\GammaO(2))=0.
\]
于是这条曲线既必须对应一个新形式，
又不可能对应任何新形式。
这就是矛盾。

\textbf{引入之后能得到什么？}
这一步完成之后，
FLT 就不再是一个外部目标，
而是整个模性理论的一条推论。
它把本书从第1章到第10章的所有主线一次性收束起来：
Fourier 系数、Jacobians、模参数化、Tate 模、Deligne 表示、水平下降与 $R=T$ 全都在同一个证明里找到位置。

\begin{figure}[ht]
\centering
\begin{tikzpicture}[>=Stealth, font=\small, node distance=1.15cm]
  \node[draw, rounded corners, fill=red!8, minimum width=5.6cm, minimum height=0.9cm] (a1) {假设 $x^n+y^n=z^n$ 有正整数解};
  \node[draw, rounded corners, fill=orange!14, below=of a1, minimum width=5.0cm, minimum height=0.9cm] (a2) {归约到 $n=4$ 或 $n=p\ge 5$};
  \node[draw, rounded corners, fill=yellow!16, below=of a2, minimum width=5.6cm, minimum height=0.9cm] (a3) {$n=4$ 被无穷递降排除；$n=p$ 导出 Frey 曲线};
  \node[draw, rounded corners, fill=green!10, below=of a3, minimum width=5.4cm, minimum height=0.9cm] (a4) {Frey 曲线半稳定，所以可用 \cref{thm:wiles}};
  \node[draw, rounded corners, fill=blue!8, below=of a4, minimum width=5.6cm, minimum height=0.9cm] (a5) {对每个 $q\mid abc$，由 \cref{thm:ribet} 逐步删去级别};
  \node[draw, rounded corners, fill=purple!10, below=of a5, minimum width=4.8cm, minimum height=0.9cm] (a6) {最终落到级别 $2$ 的权 $2$ 新形式};
  \node[draw, rounded corners, fill=gray!14, below=of a6, minimum width=4.6cm, minimum height=0.9cm] (a7) {$\Sk_2(\GammaO(2))=0$，矛盾};

  \draw[->, very thick] (a1) -- (a2);
  \draw[->, very thick] (a2) -- (a3);
  \draw[->, very thick] (a3) -- (a4);
  \draw[->, very thick] (a4) -- (a5);
  \draw[->, very thick] (a5) -- (a6);
  \draw[->, very thick] (a6) -- (a7);
\end{tikzpicture}
\caption{费马大定理证明的完整链条：假想解先被归约到奇素数指数，再经 Frey 曲线、模性与水平下降一路走向空空间。}
\label{fig:ch10-proof-chain}
\end{figure}


先做指数归约。

\begin{lemma}[把指数归约到 $4$ 或奇素数]
\label{lem:ch10-reduction}
若对 $n=4$ 与所有奇素数 $p\ge 5$，方程
\[
  x^n+y^n=z^n
\]
都没有正整数解，
则对所有 $n\ge 3$ 都没有正整数解。
\end{lemma}

\begin{proof}
若某个 $n\ge 3$ 存在正整数解，
考虑 $n$ 的素因子分解。
若 $n$ 有奇素因子 $p$，则令
\[
  X=x^{n/p},\qquad Y=y^{n/p},\qquad Z=z^{n/p},
\]
便得到
\[
  X^p+Y^p=Z^p,
\]
这与奇素数情形无解矛盾。
若 $n$ 没有奇素因子，
则 $n=2^r$ 且 $r\ge 2$。
此时令
\[
  X=x^{2^{r-2}},\qquad Y=y^{2^{r-2}},\qquad Z=z^{2^{r-2}},
\]
得到
\[
  X^4+Y^4=Z^4,
\]
从而也给出 $n=4$ 的解。
故只要排除 $4$ 与所有奇素数指数，
就排除了所有 $n\ge 3$。
\end{proof}

接着处理 Fermat 自己已经解决的指数 $4$。

\begin{theorem}[Fermat 的 $n=4$ 情形]
\label{thm:ch10-n4}
方程
\[
  x^4+y^4=z^2
\]
没有正整数解。
特别地，方程
\[
  x^4+y^4=z^4
\]
也没有正整数解。
\end{theorem}

\begin{proof}
反设存在正整数解，
并取 $z$ 最小的一组。
把公因子约去，可设 $(x,y)=1$；
交换 $x,y$ 后可设 $y$ 为偶数，$x$ 为奇数。
于是
\[
  x^4+y^4=z^2
\]
等价于 primitive Pythagorean triple
\[
  x^2,\ y^2,\ z.
\]
故存在互素整数 $m>n>0$，$m$ 奇、$n$ 偶，使得
\[
  x^2=m^2-n^2,
  \qquad
  y^2=2mn,
  \qquad
  z=m^2+n^2.
\]
由于 $2mn$ 是平方且 $(m,n)=1$，
唯一可能是
\[
  m=r^2,
  \qquad
  n=2s^2
\]
其中 $(r,s)=1$。
于是
\[
  x^2=r^4-4s^4=(r^2-2s^2)(r^2+2s^2).
\]
两因子都是正奇数，且互素；
其乘积是平方，故各自都是平方：
\[
  r^2-2s^2=u^2,
  \qquad
  r^2+2s^2=v^2
\]
其中 $0<u<v$，且 $u,v$ 仍为奇数。
于是
\[
  (v-u)(v+u)=4s^2.
\]
因为 $u,v$ 奇且 $(u,v)=1$，有
\[
  \gcd(v-u,v+u)=2.
\]
故存在互素整数 $a>b>0$ 使得
\[
  v+u=2a^2,
  \qquad
  v-u=2b^2.
\]
从而
\[
  u=a^2-b^2,
  \qquad
  s=ab.
\]
再代回
\[
  r^2=u^2+2s^2
\]
得
\[
  r^2=(a^2-b^2)^2+2a^2b^2=a^4+b^4.
\]
于是 $(a,b,r)$ 给出一个新的正整数解
\[
  a^4+b^4=r^2.
\]
但这里
\[
  r<r^2<m^2+n^2=z,
\]
所以我们从最小的 $z$ 构造出了更小的平方根，
这与最小性矛盾。
故 $x^4+y^4=z^2$ 无正整数解。

若 $x^4+y^4=z^4$ 有正整数解，
则取 $Z=z^2$ 便得到前一种方程的正整数解，
矛盾。
\end{proof}

现在进入真正的现代部分。

\begin{proposition}[Frey 曲线满足模性链条的假设]
\label{prop:ch10-frey-hypotheses}
设 $p\ge 5$，并假设存在费马方程的非平凡原始解
\[
  a^p+b^p=c^p,
  \qquad (a,b,c)=1,
  \qquad 2\mid b.
\]
令 $E=E_{a,b,c}$ 为 Frey 曲线。
则：
\begin{itemize}
  \item[(i)] $E$ 半稳定，因此可对 $E$ 应用 \cref{thm:wiles}；
  \item[(ii)] 对每个奇素数 $q\mid abc$，$\rhobar_{E,p}$ 在 $q$ 处无分歧；
  \item[(iii)] $\rhobar_{E,p}$ 是奇且绝对不可约，因而满足 \cref{thm:ribet} 的假设。
\end{itemize}
\end{proposition}

\begin{proof}
(i) 是 \cref{prop:ch10-frey-conductor} 的结论。
(ii) 是 \cref{prop:ch10-frey-unramified} 的结论。

对 (iii)，奇性来自行列式：
\[
  \det \rhobar_{E,p}=\chi_p,
\]
其中 $\chi_p$ 是 mod $p$ 分圆特征。
复共轭在 $E[p]$ 上的行列式等于 $-1$，
故表示是奇的。

绝对不可约性是 Frey 曲线理论中的标准事实。
若 $\rhobar_{E,p}$ 可约，
则 $E$ 将带有有理 $p$-isogeny；
结合 Frey 曲线的半稳定性、全有理 $2$-torsion 以及其特殊的 $j$-不变量形状，
可由 Mazur 的有理同源分类与 Serre 的不可约性判据推出矛盾。
因此 $\rhobar_{E,p}$ 绝对不可约。
在本章层次上，我们把这一步视为输入性的标准结论。
\end{proof}

于是可以把主定理写成一条极短但极有力的证明。

\begin{theorem}[费马大定理]
\label{thm:flt}
对任意整数 $n\ge 3$，方程
\[
  x^n+y^n=z^n
\]
没有正整数解。
\end{theorem}

\begin{proof}
由 \cref{lem:ch10-reduction}，只需处理 $n=4$ 与 $n=p\ge 5$ 为奇素数两种情形。
$n=4$ 已由 \cref{thm:ch10-n4} 解决。

现设 $p\ge 5$ 为奇素数，
反设存在正整数解
\[
  x^p+y^p=z^p.
\]
约去公因子后可设 $(x,y,z)=1$；
交换 $x,y$ 后可设 $2\mid y$。
记
\[
  a=x,\qquad b=y,\qquad c=z,
\]
并构造 Frey 曲线
\[
  E=E_{a,b,c}\colon y^2=x(x-a^p)(x+b^p).
\]

由 \cref{prop:ch10-frey-hypotheses}(i)，
$E$ 是半稳定的。
因此可应用 \cref{thm:wiles}：
存在某个权 $2$ 的归一化特征新形式
\[
  f\in \Sk_2(\GammaO(N(E)))
\]
使得 $E$ 模，等价地，其 $p$-进表示与 $f$ 的表示一致。
特别地，mod $p$ 约化
\[
  \rhobar_{E,p}
\]
来自某个级别 $N(E)=2\,\mathrm{rad}_{\mathrm{odd}}(abc)$ 的权 $2$ 新形式。

接下来对每个奇素数 $q\mid abc$ 应用 \cref{thm:ribet}。
由 \cref{prop:ch10-frey-hypotheses}(ii)，
$\rhobar_{E,p}$ 在这些 $q$ 处无分歧；
由 \cref{prop:ch10-frey-hypotheses}(iii)，
它满足绝对不可约假设。
因此每应用一次 \cref{thm:ribet}，
就能把级别中的一个奇素数删去。
对所有 $q\mid abc$ 反复进行，
最终得到：
$\rhobar_{E,p}$ 来自某个级别 $2$ 的权 $2$ 新形式。

但第5章已经用亏格公式算出
\[
  \dim \Sk_2(\GammaO(2))=g(X_0(2))=0.
\]
因此根本不存在级别 $2$ 的权 $2$ 新形式。
这与上一步得到的存在性矛盾。
故奇素数指数情形也不可能有正整数解。

结合 $n=4$ 的情形，
便知对所有 $n\ge 3$ 都无正整数解。
\end{proof}

\textbf{注。}
这段证明最值得反复体会的地方，
不是“矛盾终于出现了”，
而是矛盾出现的位置异常干净：
它既不来自巨大计算，
也不来自偶然同余，
而是来自一个空的模形式空间。
Fermat 的方程若有解，
就会迫使空空间里长出一个新形式；
而空空间的沉默，
就是全书最后的判决。
